% 7点图：贯穿四点图 {12,13,23,24,34} 加边 {5,6} 与孤立点 7。外层按编号扫描、每次启动BFS。
% 分量 {1,2,3,4}、{5,6}、{7}；父边 {1,2},{1,3},{2,4},{5,6}，共 4 = 7 - 3 条。
\begin{tikzpicture}[font=\figurefont\fontsize{8}{10}\selectfont,
  vtx/.style={circle,draw=FigureStroke,line width=.8pt,minimum size=6.5mm,inner sep=0pt,text=FigureInk},
  root/.style={double=white,double distance=.6pt,line width=.6pt},
  tree/.style={draw=FigureBlue,line width=1.4pt},
  other/.style={draw=FigureMuted,line width=.7pt,dashed},
  box/.style={draw=FigureGrid,line width=.6pt,rounded corners=1.8pt},
  note/.style={font=\FigureNoteFont\fontsize{7.5}{9}\selectfont,text=FigureMuted}]
  \draw[box] (-0.5,-1.3) rectangle (3.1,1.3);
  \draw[box] (3.5,-1.3) rectangle (5.1,1.3);
  \draw[box] (5.5,-1.3) rectangle (6.6,1.3);
  \node[vtx,fill=FigureBlueFill,root] (v1) at (0,0) {1};
  \node[vtx,fill=FigureBlueFill] (v2) at (1.3,0.8) {2};
  \node[vtx,fill=FigureBlueFill] (v3) at (1.3,-0.8) {3};
  \node[vtx,fill=FigureBlueFill] (v4) at (2.6,0) {4};
  \node[vtx,fill=FigureGreenFill,root] (v5) at (4.3,0.6) {5};
  \node[vtx,fill=FigureGreenFill] (v6) at (4.3,-0.6) {6};
  \node[vtx,fill=FigureYellowFill,root] (v7) at (6.05,0) {7};
  \draw[other] (v2)--(v3) (v3)--(v4);
  \draw[tree] (v1)--(v2) (v1)--(v3) (v2)--(v4) (v5)--(v6);
  \node at (1.3,-1.6) {编号 1};
  \node at (4.3,-1.6) {编号 2};
  \node at (6.05,-1.6) {编号 3};
  \node[note,anchor=west] at (6.95,0.45) {双圈：外层扫描的起点};
  \node[note,anchor=west] at (6.95,0) {蓝色粗线：发现父边};
  \node[note,anchor=west] at (6.95,-0.45) {虚线：未选入森林的边};
\end{tikzpicture}
